\documentclass[../main.tex]{subfiles}
\providecommand{\mainx}{..}
\begin{document}
\section{Algebra of sets} 
\label{sec:setalgebra}
A \emph{set}\index{set} is an unordered collection of distinct elements.
If we know the elements in a set, we may denote the set by these elements, e.g., $\{a,c,b\}$ denotes a set whose members are exactly $a$, $b$, and $c$.

Two sets of particular importance are the empty set, denoted by $\EmptySet$, which has no members, and the \emph{universal set}, in which every element of interest is a member.

A \emph{finite set} has a finite number of elements.
For example, $\{ 1, 3, 5 \}$ is a finite set with three elements.
When sets $\Set{A}$ and $\Set{B}$ are \emph{isomorphic}, denoted by $\Set{A} 
\cong \Set{B}$, they can be put into a one-to-one correspondence (bijection), e.g., $\{b,a,c\} \cong \{1,2,3\}$.

The cardinality of a finite set $\Set{A}$ is the number of elements in the set, denoted by $\Card{\Set{A}}$, e.g., $\Card{\left\{ 1, 3, 5\right\}} = 3$.
A \emph{countably infinite set} is isomorphic to the set of \emph{natural numbers} $\NatSet \coloneqq \{1,2,3,4,5,\ldots\}$.


Given two elements $a$ and $b$, an ordered pair of $a$ then $b$ is denoted by $\Pair{a}{b}$, where $\Pair{a}{b} = \Pair{c}{d}$ \emph{if and only if} $a = c$ and $b = d$.
Ordered pairs are non-commutative and non-associative, i.e., $\Pair{a}{b} \neq \Pair{b}{a}$ if $a \neq b$ and $\Pair{a}{\Pair{b}{c}} \neq \Pair{\Pair{b}{a}}{c}$.

Related to the ordered pair is the Cartesian product.
\begin{definition}
	The set $\Set{X} \times \Set{Y} \coloneqq \left\{\Pair{x}{y} \colon x \in 
	\Set{X} \land y \in \Set{Y}\right\}$ is the Cartesian product of sets $\Set{X}$ 
	and $\Set{Y}$.
\end{definition}
By the non-commutative and non-associative property of ordered pairs, the Cartesian product is non-commutative and non-associative.
However, they are isomorphic, i.e., $\Set{X} \times \Set{Y} \cong \Set{Y} \times \Set{X}$.

A \emph{tuple} is a generalization of order pairs which can consist of an arbitrary number of elements, e.g., $\Tuple{x_1,x_2,\ldots,x_n}$.
\begin{definition}[$n$-fold Cartesian product]
	The $n$-ary Cartesian product of sets $\Set{X}[1],\ldots,\Set{X}[n]$, is given by $\Set{X}[1] \times \cdots \times \Set{X}[n] = \left\{\Tuple{x_1,\ldots,x_n} \colon x_1 \in \Set{X}[1] \land \cdots \land x_n \in \Set{X}[n]\right\}$.
\end{definition}
Note that
	$\Set{X}[1] \times \Set{X}[2] \times \Set{X}[3] \cong \Set{X}[1] \times \left(\Set{X}[2] \times \Set{X}[3]\right) \cong \left(\Set{X}[1] \times \Set{X}[2]\right) \times \Set{X}[3]$,
thus we may implicitly convert between them without ambiguity.

If each set in the $n$-ary Cartesian product is the same, the power notation may be used, e.g., $\Set{X}^3 \coloneqq \Set{X} \times \Set{X} \times \Set{X}$.
As special cases, $\Set{X}^0 \coloneqq \{ \EmptySet \}$ and $\Set{X}^1 \coloneqq \Set{X}$.

A \emph{binary relation} over sets $\Set{A}$ and $\Set{B}$ is any subset of $\Set{A} \times \Set{B}$.
A fundamental relation is the member-of relation, where $\SetContains[x][\Set{A}]$ denotes that an object $x$ is a member of a set $\Set{A}$.
A set $\Set{A}$ is a \emph{subset} of a set $\Set{B}$ if every member of $\Set{A}$ is a member $\Set{B}$, denoted by $\Set{A} \subseteq \Set{B}$.
The subset relation forms a \emph{partial order}, i.e., if $\Set{A} \subseteq \Set{B}$ and $\Set{B} \subseteq \Set{C}$ then $\Set{A} \subseteq \Set{C}$ and 
if $\Set{A} \subseteq \Set{B}$ and $\Set{B} \subseteq \Set{A}$ then $\Set{A}$ and $\Set{B}$ are \emph{equal}, denoted by $\Set{A} = \Set{B}$.


\begin{definition}
	Set builder notation
\end{definition}


\begin{definition}
	A \emph{function} of type $\Set{X} \mapsto \Set{Y}$ is a binary relation on $\Set{X} \times \Set{Y}$ with the constraint that each $x \in \Set{X}$ is paired with exactly one $y \in \Set{Y}$.
\end{definition}
A function of type $\Set{X} \mapsto \Set{Y}$ has a domain $\Set{X}$ and a codomain $\Set{Y}$.
Since every $x \in \Set{X}$, given a pair $\Pair{x}{y} \in \Fun{f}$, $y$ may also be denoted by $\Fun{f}(x)$.

The \emph{power set} of a set $\Set{A}$, denoted by $\PS{\Set{A}}$, is the set of sets that contains all of the possible subsets of $\Set{A}$, e.g., $\PS{\{a, b\}}= \left\{ \EmptySet, \{a\}, \{b\}, \{a, b\} \right\}$.

A predicate is a function that maps elements in its domain to true (denoted by $1$) or false (denoted by $0$).
A predicate function of particular importance is the indicator function
\begin{equation}
	\SetIndicator{\Set{A}} \colon \Set{X} \mapsto \{0,1\}
\end{equation}
defined as
\begin{equation}
	\SetIndicator{\Set{A}}(x) \coloneqq
	\begin{cases}
		0 & \text{if $x \notin \Set{A}$}\,,\\
		1 & \text{if $x \in \Set{A}$}\,.
	\end{cases}
\end{equation}

The indicator function admits the construction of predicates for any relation, e.g., a binary predicate $\operatorname{P}$ for a binary relation $\Set{R} \subseteq \Set{A} \times \Set{B}$ is defined as $\operatorname{P}(x_1,x_2) \coloneqq \SetIndicator{\Set{R}}(\Pair{x_1}{x_2})$.
Denoting the \emph{universal set} by $\Set{X}$, all the relations mentioned previously are \emph{binary predicates}, such as $\SetContains \colon \Set{X} \times \PS{\Set{X}} \mapsto \{0,1\}$ and $\subseteq \colon \PS{\Set{X}} \times \PS{\Set{X}} \mapsto \{0,1\}$.

A few important operations on sets are \emph{set-union} $\SetUnion \colon \PS{\Set{X}} \times \PS{\Set{X}} \mapsto \PS{\Set{X}}$, \emph{set-intersection}, $\SetIntersection \colon \PS{\Set{X}} \times \PS{\Set{X}} \mapsto \PS{\Set{X}}$, and \emph{set-complement} $\SetComplement \colon \PS{\Set{X}} \mapsto \PS{\Set{X}}$, respectively defined as
\begin{align}
	\SetUnion[\Set{A}][\Set{B}]			&\coloneqq \SetBuilder{ x \in \Set{X}}{x \in \Set{A} \lor x \in \Set{B}}\,,\\
	\SetIntersection[\Set{A}][\Set{B}]	&\coloneqq \SetBuilder{ x \in \Set{X}}{x \in \Set{A} \land x \in \Set{B}}\,,\,\text{and}\\
	\SetComplement[\Set{A}]				&\coloneqq \SetBuilder{ x \in \Set{X}}{x \notin \Set{A}}\,.
\end{align}
where $\lor$ and $\land$ are respectively the logical-connectives \emph{or} and \emph{and}.
If $\SetIntersection[\Set{A}][\Set{B}] = \EmptySet$, then we say $\Set{A}$ and $\Set{B}$ are \emph{disjoint} sets.

\subsection{Approximate sets}
\label{sec:asets}
Given an objective set $\Set{A}$, any element that is a member of $\Set{A}$ is denoted a \emph{positive} (of $\Set{A})$ and otherwise the element is denoted a \emph{negative}.
Suppose $\hat{\Set{A}}$ is used as an \emph{approximation} of $\Set{A}$.
If the \emph{only} information we have about $\Set{A}$ is given by $\hat{\Set{A}}$, then we may perform membership tests on $\hat{\Set{A}}$ to make \emph{predictions} or \emph{estimations} about $\Set{A}$.

There are two ways a binary prediction can be false.
\begin{enumerate}
	\item A \emph{false positive} occurs if a negative of the objective set is predicted to be a positive. False positives are also known as \emph{type I errors}.
	The complement of false positives are \emph{true negatives}.
	\item A \emph{false negative} occurs if a positive of the objective set is predicted to be a negative. False negatives are also known as \emph{type II errors}.
	The complement of false negatives are \emph{true positives}.
\end{enumerate} 

If we denote the set of false positives by $\PlainSet{FP}$, true positives by $\PlainSet{TP}$, false negatives by $\PlainSet{FN}$, and true negatives by $\PlainSet{TN}$, then the objective set $\Set{A}=\SetUnion[\PlainSet{FN}][\Set{TP}]$ and the approximate set $\hat{\Set{A}}=\SetUnion[\PlainSet{TP}][\PlainSet{FP}]$.
See \cref{fig:ex_approx_set} for an illustration.
\begin{figure}[ht]
	\caption{An approximate set $\hat{\Set{A}}$ of an objective set $\Set{A}$}
	\label{fig:ex_approx_set}
	\centering
	\def\svgwidth{\columnwidth/4}
	%% Creator: Inkscape inkscape 0.92.1, www.inkscape.org
%% PDF/EPS/PS + LaTeX output extension by Johan Engelen, 2010
%% Accompanies image file 'aset_fp_fn.pdf' (pdf, eps, ps)
%%
%% To include the image in your LaTeX document, write
%%   \input{<filename>.pdf_tex}
%%  instead of
%%   \includegraphics{<filename>.pdf}
%% To scale the image, write
%%   \def\svgwidth{<desired width>}
%%   \input{<filename>.pdf_tex}
%%  instead of
%%   \includegraphics[width=<desired width>]{<filename>.pdf}
%%
%% Images with a different path to the parent latex file can
%% be accessed with the `import' package (which may need to be
%% installed) using
%%   \usepackage{import}
%% in the preamble, and then including the image with
%%   \import{<path to file>}{<filename>.pdf_tex}
%% Alternatively, one can specify
%%   \graphicspath{{<path to file>/}}
%% 
%% For more information, please see info/svg-inkscape on CTAN:
%%   http://tug.ctan.org/tex-archive/info/svg-inkscape
%%
\begingroup%
  \makeatletter%
  \providecommand\color[2][]{%
    \errmessage{(Inkscape) Color is used for the text in Inkscape, but the package 'color.sty' is not loaded}%
    \renewcommand\color[2][]{}%
  }%
  \providecommand\transparent[1]{%
    \errmessage{(Inkscape) Transparency is used (non-zero) for the text in Inkscape, but the package 'transparent.sty' is not loaded}%
    \renewcommand\transparent[1]{}%
  }%
  \providecommand\rotatebox[2]{#2}%
  \ifx\svgwidth\undefined%
    \setlength{\unitlength}{437.78998862bp}%
    \ifx\svgscale\undefined%
      \relax%
    \else%
      \setlength{\unitlength}{\unitlength * \real{\svgscale}}%
    \fi%
  \else%
    \setlength{\unitlength}{\svgwidth}%
  \fi%
  \global\let\svgwidth\undefined%
  \global\let\svgscale\undefined%
  \makeatother%
  \begin{picture}(1,1.05361373)%
    \put(-0.00881474,0.603155){\color[rgb]{0,0,0}\makebox(0,0)[lb]{\smash{}}}%
    \put(-0.01615682,0.61783922){\color[rgb]{0,0,0}\makebox(0,0)[lb]{\smash{}}}%
    \put(0,0){\includegraphics[width=\unitlength,page=1]{aset_fp_fn.pdf}}%
    \put(0.78,0.06){\color[rgb]{0,0,0}\makebox(0,0)[lb]{\smash{$\Set{A} = \SetUnion[\PlainSet{TP}][\PlainSet{FN}]$}}}%
    \put(0.01737703,1.02584642){\color[rgb]{0,0,0}\makebox(0,0)[lb]{\smash{$\hat{\Set{A}} = \SetUnion[\PlainSet{TP}][\PlainSet{FP}]$}}}%
    \put(0.67844448,0.84067833){\color[rgb]{0,0,0}\makebox(0,0)[lb]{\smash{$\PlainSet{FP}$}}}%
    \put(0.84976816,0.43053962){\color[rgb]{0,0,0}\makebox(0,0)[lb]{\smash{$\PlainSet{FN}$}}}%
    \put(0.50885122,0.45649776){\color[rgb]{0,0,0}\makebox(0,0)[lb]{\smash{$\PlainSet{TP}$}}}%
    \put(0.03987408,0.45476734){\color[rgb]{0,0,0}\makebox(0,0)[lb]{\smash{$\PlainSet{TN}$}}}%
    \put(0,0){\includegraphics[width=\unitlength,page=2]{aset_fp_fn.pdf}}%
  \end{picture}%
\endgroup%

\end{figure}

If we only have access to the approximation $\hat{\Set{A}}$, we do not have the ability to partition the universe into the disjoint sets $\PlainSet{FP}$, $\PlainSet{TP}$, $\PlainSet{FN}$, and $\PlainSet{TN}$ as demonstrated in \cref{fig:ex_approx_set}.
However, we can quantify the degree of \emph{uncertainty} about the elements that it predicts to be positive or negative.
The false positive and true negative rates are given by the following.
\begin{definition}
	\label{def:fprate}
	The \emph{false positive rate} is the proportion of predictions that are \emph{false positives} as given by
	\begin{equation}
	\hat{\fprate} = \frac{\Card{\PlainSet{FP}}}{\Card{\PlainSet{FP}} + \Card{\PlainSet{TN}}}
	\end{equation}
	and the \emph{true negative rate} is given by $\hat{\tnrate} = 1 - \hat{\fprate}$.
\end{definition}
The true positive and false negative rates are given by the following.
\begin{definition}
	The \emph{true positive rate} is the proportion of predictions that are \emph{true positives} as given by
	\begin{equation}
	\hat{\tprate} = \frac{\Card{\PlainSet{TP}}}{\Card{\PlainSet{TP}} + \Card{\PlainSet{FN}}}
	\end{equation}
	and the \emph{false negative rate} is given by $\hat{\fnrate} = 1 - \hat{\tprate}$.
\end{definition}

The \emph{probabilities} of the four possible predictive outcomes are given by 
\cref{tbl:contingency_table}.
\begin{table}[ht]
	\centering
	\begin{tabular}{@{} l l l @{}}
		\toprule
		& \textbf{positive} & \textbf{negative}\\
		\midrule
		\textbf{predict positive} & $\hat{\tprate}=1-\hat{\fnrate}$ & 
		$\hat{\fprate}=1-\hat{\tnrate}$\\
		\textbf{predict negative} & $\hat{\fnrate}=1-\hat{\tprate}$ & 
		$\hat{\tnrate}=1-\hat{\fprate}$\\
		\bottomrule
	\end{tabular}
	\caption{The $2 \times 2$ contingency table of outcomes for approximate sets.}
	\label{tbl:contingency_table}        
\end{table}

\end{document}